\documentclass[11pt,a4paper]{article}
\usepackage[T1]{fontenc}
\usepackage{lmodern,amsmath,amssymb,microtype,booktabs}
\usepackage[margin=25mm]{geometry}
\usepackage[colorlinks=true,urlcolor=blue,linkcolor=blue]{hyperref}
\newcommand{\Om}{\Omega}
\title{Exponential amplification in the dS$_3$ crossing kernel:\\
full-factor tail bounds and an obstruction to direct summation}
\author{Alan Day \and Chengheng\thanks{AI research assistant (OpenAI).}}
\date{Working research note --- 31 October 2026}
\begin{document}\maketitle
\begin{abstract}
An accurate approximation to one crossing factor need not give an accurate
physical matrix element. We construct exact rational polynomial certificates
for the tails of all four Omega factors in the known scalar LPRS formula,
including the digamma series at a removable collision. We propagate the
bounds through the two shadow terms and derive an elementary expression
for their coefficient moduli. One coefficient grows as
$\exp(2\pi M)$ in the specified heavy-observer limit.
Direct partial sums can therefore fail by many orders of magnitude even
when their arithmetic is stable and their individual absolute tails are small.
This identifies an analytic obstacle; it does not establish a heavy-mass
distributional limit or a de Sitter holographic dual.
\end{abstract}

\section{Scope and a general series certificate}
We retain radius one, $G_N=0$, proper $SO^+(3,1)$, BD scalar observers
with a fixed positive mass-label lower bound, KG-normalized matter,
and the previously defined half-period regulated local creation blocks.
The object is the scalar-to-scalar element $U_{00}$ of the full
Racah transform. The scalar subblock is not asserted to be unitary.
External labels are $(a,q,c,p)$ and internal labels $(b,d)$, all positive.
The chiral weights obey $h_i=(1+i\nu_i)/2$ with $\bar h_i=h_i$,
not complex conjugation.

Every ordinary hypergeometric series needed in LPRS (3.38) has terms
\[
 t_n=\frac{\prod_{j=1}^4(\frac12+i\alpha_j)_n}
 {n!\prod_{k=1}^3(1+i\beta_k)_n},\qquad
 \alpha_j,\beta_k\in\mathbb R,\quad t_0=1.
\]
The definition is standard \cite{dlmf16}. Set
\[
 A(x)=\prod_j[(x+\tfrac12)^2+\alpha_j^2],\quad
 B(x)=(x+1)^2\prod_k[(x+1)^2+\beta_k^2],
\]
\[
 P(x)=B(x)(x+1)^3-A(x)(x+2)^3.
\]
The degree-eleven terms cancel; $P$ has degree ten and leading coefficient
one, independent of the parameters. Choose an integer $N_0$ for which
all coefficients of $P(N_0+y)$ are nonnegative. This implies
\begin{equation}
 \left|\frac{t_{n+1}}{t_n}\right|
 \le\left(\frac{n+1}{n+2}\right)^{3/2},\qquad n\ge N_0 .
 \label{ratio}
\end{equation}
Such a shift exists for every fixed parameter set, since the derivatives
of a positive-leading polynomial are eventually positive.
For rational imaginary parts we check the coefficients with exact rational
arithmetic. This is a sufficient certificate, not a minimal cutoff
or a bound uniform in growing masses.

\newpage
\section{Tails at and away from a collision}
For $N\ge N_0$, \eqref{ratio} implies
$|t_n|\le |t_N|((N+1)/(n+1))^{3/2}$ for $n\ge N$.
The integral test gives the exact-arithmetic tail bound
\begin{equation}
 \left|\sum_{n=N}^{\infty}t_n\right|
 \le |t_N|(2N+3).
 \label{ordinary}
\end{equation}
Unlike an asymptotic zeta acceleration, no discarded inverse-power
remainder is involved. The price can be a very loose bound.

At an exact collision let $H$ be either $h$ or $1-h$,
$h'_0=H+h_1+h_3-1$, and
\[
 (A_1,A_2,A_3,A_4)
 =(H+h_1+h_2-1,H+h_3+h_4-1,H+h_3-h_4,H+h_1-h_2),
\]
\[
 (B_1,B_2,B_3,B_4)=(2h'_0,2H,1,1).
\]
The preceding merged-series derivation gives
\[
 \Om_0=\sum_{n=0}^{\infty}t_n q_n,\qquad
 q_n=\sum_{j=1}^4[\psi(n+B_j)-\psi(n+A_j)],
\]
where the $t_n$ denominator is $(2h'_0)_n(2H)_n(n!)^2$.
All $A_j$ have real part $1/2$ and all $B_j$ real part one.
Using the trigamma series \cite{dlmf5} along the straight segment between
the two arguments,
\[
 |q_n|\le L\left[\frac1{n+1/2}+\frac1{(n+1/2)^2}\right],
 \qquad L=\sum_j|B_j-A_j|.
\]
Define $C_N=L(1+(N+1/2)^{-1})(N+1)/(N+1/2)$.
A second integral test yields
\begin{equation}
 \left|\sum_{n=N}^{\infty}t_nq_n\right|
 \le |t_N|C_N\left[\frac1{N+1}+\frac23\right].
 \label{weighted}
\end{equation}
The implementation advances $q_n$ using the digamma recurrence, retaining
all complex phases. Away from collision the two ordinary series receive
their exact Gamma coefficients, and their two bounds are added in modulus.
That bound need not remain useful as the collision is approached.

Write the complete physical element as
\[
 U_{00}=C_+\Om(h,h')\Om(1-h,h')
       +C_-\Om(h,1-h')\Om(1-h,1-h').
\]
For approximate factors $x_\sigma,y_\sigma$ with tail bounds
$e_{x,\sigma},e_{y,\sigma}$, the propagated bound is
\begin{equation}
 |U_{00}-\widehat U_{00}|
 \le\sum_{\sigma=\pm}|C_\sigma|
 (|x_\sigma|e_{y,\sigma}+|y_\sigma|e_{x,\sigma}
                         +e_{x,\sigma}e_{y,\sigma}).
 \label{full}
\end{equation}
Equations \eqref{ordinary}--\eqref{full} concern exact finite sums and
exact Gamma coefficients. Evaluating their right-hand sides with
mpmath is not an interval certificate for rounding or special functions.

\newpage
\section{The exponentially amplified shadow coefficient}
The coefficients include the physical phase normalization of our earlier
crossing dictionary, whose modulus is $bd/\pi^4$.
LPRS (3.35)--(3.38) \cite{lprs} then imply an elementary modulus formula.
Define
\[
 x_1=\frac{b+a-q}{2},\quad x_2=\frac{b-a+q}{2},\quad
 x_3=\frac{-b+p-c}{2},\quad x_4=\frac{-b+c-p}{2}.
\]
The Gamma reflection identities give
\[
 |\Gamma(\tfrac12+ix)|^2=\frac{\pi}{\cosh\pi x},\qquad
 |\Gamma(1+ib)|^2=\frac{\pi b}{\sinh\pi b}.
\]
The four-Gamma common prefactor has modulus
$\pi\sinh(\pi b)/(4b\sqrt{\prod_j\cosh(\pi x_j)})$.
Each shadow coefficient has modulus $\pi/d$; the remaining Gamma
ratios have unit modulus because their half-integer arguments pair
with opposite imaginary parts. Finally,
$|\sin(-\pi/2+iy)|=\cosh y$. Consequently,
\begin{equation}
 |C_\sigma|=
 \frac{\sinh(\pi b)}
 {4\pi^2\sqrt{\prod_j\cosh(\pi x_j)}}
 \cosh\frac{\pi(\sigma d-a-p)}2
 \cosh\frac{\pi(\sigma d-q-c)}2 .
 \label{coefficient}
\end{equation}
This is an exact modulus identity, not a phase convention or a replacement
of the full complex kernel by positive quantities.

Now fix $u,v,p,q$ and take
$c=M$, $b=M+u$, $a=M+u+v$, $d=M+v$. Then
\begin{align}
 \frac{|C_-|}{|C_+|}
 &=\frac{\exp(2\pi M+\frac\pi2(u+3v+p+q))}
 {4\cosh\frac{\pi(u+p)}2\cosh\frac{\pi(v-q)}2}
 \nonumber\\[-2pt]
 &\quad\times(1+e^{-\pi(d+a+p)})(1+e^{-\pi(d+q+c)}).
 \label{amplification}
\end{align}
Moreover,
\[
 |C_+|\longrightarrow
 \frac{e^{\pi(u-v+p+q)/4}}{4\pi^2}
 \sqrt{\cosh\frac{\pi(u+p)}2\cosh\frac{\pi(v-q)}2}>0.
\]
Thus $|C_-|$ grows as $e^{2\pi M}$ with a nonzero fixed coefficient.
An independent absolute-error estimate for
$\Om(h,1-h')$ must offset this exponential gain whenever its partner
is not comparably small. Polynomial accuracy of a different Omega
cannot supply that control.

This does not prove that the physical matrix element diverges.
It explains why an exponentially small factor or cancellations must be
controlled before taking the heavy-mass limit of this representation.

\newpage
\section{Numerical evidence, limitations, and next calculation}
We use $p=q=0.6$, $u=0.3$, $v=-0.2$ on the collision surface.
All tested polynomial certificates allow $N=64$; we also sum $N=128$.
The latter sums at 110 and 150 decimal digits agree relatively to
better than $10^{-106}$ for the two larger masses. They are nevertheless
poor approximations to the complete physical element:
\begin{center}
\begin{tabular}{rrrr}
\toprule
$M$ & $N$ & observed $|\widehat U-U_{\rm ref}|$ & full tail bound\\
\midrule
2 &128& $2.862\,10^{-3}$ & $3.589\,10^{-1}$\\
8 &128& $1.064\,10^{7}$ & $3.391\,10^{8}$\\
32&128& $4.982\,10^{43}$ & $7.527\,10^{44}$\\
\bottomrule
\end{tabular}
\end{center}
The references are the previously saved complete-kernel calculations.
The $M=2$ reference retains its known prototype bias; no reference is
being treated as an exact interval enclosure. These observed errors are
far larger than that bias. All nine runs, including the failed
approximations, are preserved.

At $M=32$, the modulus ratio in \eqref{amplification} is approximately
$10^{86.7155}$. The computed tail of the first Omega in the minus-shadow
term is only about $4.73\,10^{-42}$, but its amplification is enormous.
Holding the other computed factor fixed, an absolute $10^{-6}$ budget
for this product requires an error near $6.28\,10^{-93}$ in that Omega.
This is a sensitivity calculation at the stored finite sum, not a
certificate that the exact partner or complete product has that error.

Independent Fraction arithmetic reconstructs every shifted polynomial;
Decimal arithmetic replays each first-omitted modulus, ordinary or
weighted tail, full-product propagation, and complex reference residual.
All 407 checks pass. Passing these checks verifies the bound machinery
and the retained failures, not the accuracy of these direct sums.
The exact coefficient identity is independently checked against the
Gamma expression at 150 digits. Versions and source/input hashes
accompany the saved results.

The preceding collision note's separate 100-digit hypergeometric
parameter-derivative validation remains running at this checkpoint.
It is neither required for the rational certificates here nor counted
as a passed check. Both stages remain unfrozen pending final project
integration.

The next analytic problem is the exponentially small
$\Om(h,1-h')$ on this physical slice, with its phase and multiplying
coefficient kept together. Candidates are an Euler-contour deformation
or a complete-kernel integral in which the large coefficient has
already been combined. A fixed-label tail certificate alone cannot
justify an exchange with the physical mass integrals. The full crossed
heavy-observer limit, Bose free limit, full Fock domain, finite-$G_N$
dynamics, and microscopic holographic dictionary remain open.

\begin{thebibliography}{9}
\bibitem{lprs} J. Liu, E. Perlmutter, V. Rosenhaus and D. Simmons-Duffin,
\emph{$d$-dimensional SYK, AdS Loops, and $6j$ Symbols},
\href{https://arxiv.org/html/1808.00612v3}{arXiv:1808.00612v3},
Eqs.\ (3.35)--(3.38).
\bibitem{dlmf16} NIST DLMF,
\href{https://dlmf.nist.gov/16.2}{Sec.\ 16.2}, generalized hypergeometric series.
\bibitem{dlmf5} NIST DLMF,
\href{https://dlmf.nist.gov/5.15}{Sec.\ 5.15}, trigamma series;
\href{https://dlmf.nist.gov/5.4}{Sec.\ 5.4}, Gamma modulus identities.
Consulted 31 October 2026.
\end{thebibliography}
\end{document}
